\section{Datasets}

This section describes the datasets employed in this study and the augmentation methods applied during training. Two datasets were used: one collected independently in Japan (FruitsPark) and a publicly available urban tree dataset (UrbanStreetTree).

\subsection{FruitsPark Dataset}
The FruitsPark dataset was independently collected at a public fruit park in Japan between June 2024 and May 2025 as part of a field survey conducted in a designated experimental area. It consists of images from 10 tree classes. Images were captured using an iPhone 15 Pro and a Pixel 7 Pro, unified to a resolution of 1080 $\times$ 1920 pixels. The dataset covers all four seasons and includes diverse backgrounds such as soil, other trees, playground equipment, and signboards within the park. To prevent data leakage between training and validation, samples were divided by individual trees and collection periods. Sample images from the dataset are shown in Figure~5. Hereafter, this dataset is referred to as FruitsPark.

Cherry\quad Kiwi\quad Peach\quad Persimmon\quad Pomegranate

Figure 5. Samples of FruitsPark dataset

\subsection{UrbanStreetTree Dataset}
The UrbanStreetTree dataset (Branch subset) consists of urban street tree images collected across ten cities in China. \cite{ref11} Images were captured using multiple devices at a resolution of 3008 $\times$ 3970 pixels. Since the dataset includes mid-range images containing people, vehicles, and buildings in the background, it is well suited for evaluating classification performance in complex urban environments. Sample images from the dataset are shown in Figure~6. Hereafter, this dataset is referred to as UrbanStreetTree.

Magnolia liliflora Desr\quad Ginkgo biloba\quad Platanus\quad Albizia julibrissin\quad Legerstroemia indica

Figure 6. Samples of UrbanStreetTree dataset \cite{ref11}

\section{Data Augmentation}
\subsection{Standard Photometric Augmentation}
During training, standard photometric augmentations were applied, including color, brightness, and blur variations. Specifically, random brightness and contrast adjustments ($\pm 0.2$), hue–saturation–value transformations, motion blur (3–5 pixels), resizing to $224 \times 224$, and normalization (ImageNet mean and variance) were performed. During validation, only resizing and normalization were applied. These operations improved model robustness against variations in lighting and capture conditions.

\subsection{Depth-based Background Replacement}
In addition to standard photometric augmentation, a depth-guided augmentation using relative depth maps estimated  was probabilistically applied during training. Foreground regions (trees) were preserved, while background regions were replaced or degraded using external natural images. Morphological operations and feathering were used to blend the boundaries naturally. This augmentation aims to suppress background-dependent overfitting and encourage attention to intrinsic structural features such as trunks and branches. The method was not applied during validation or inference. The detailed configuration is summarized in Table~1, the overall procedure is illustrated in Figure~7.

Aalbizia\quad Ginkgo\quad Lagerst-roemia\quad Salix\quad Zelkova

Figure 7. Sample images generated by depth-guided augmentation.

\subsection{Dataset Comparison}
Table 2 and Table 3 summarize the datasets used in this study: the FruitsPark dataset (10 classes, 2,483 images) and the UrbanStreetTree dataset (13 classes, 4,966 images) \cite{ref11}. 

Table 1. Configuration of depth-guided background replacement augmentation

Processing Step\quad Description / Setting\quad Probability

Depth estimation and foreground extraction\quad Depth estimated using MiDaS (DPT\_Hybrid model); foreground (tree) region extracted and only the largest connected component retained\quad 100\%

Background replacement\quad Background region replaced with natural scenes from the Places365 dataset (depth matching enabled)\quad 80\%

Mask processing\quad Erosion: 5 px; Dilation: 3 px; Feathering: 13 px\quad 100\%

Blur effect\quad Gaussian blur applied with strength 3–10\quad 90\%

Color transformation\quad Brightness, contrast, and saturation scaled by 0.65–1.35; hue shifted by $\pm 0.10$\quad 55\%

JPEG compression\quad JPEG compression applied with quality range of 16–50\quad 70\%

Table 2. Data splits before and after augmentation

Dataset\quad Train\quad Val\quad Test

FruitsPark (original)\quad 580\quad 210\quad 0

FruitsPark (augmented)\quad 2,273\quad 210\quad 0

UrbanStreetTree (original)\quad 1,193\quad 149\quad 143

UrbanStreetTree (augmented)\quad 4,674\quad 149\quad 143

Table 3. Comparison of the datasets used in this study

Item\quad FruitsPark\quad UrbanStreetTree

Number of classes\quad 10\quad 13

Number of images\quad 790\quad 1,485

Number of images (augmented)\quad 2,483\quad 4,966

Data split\quad Training, validation (no test set)\quad Training, validation, test

Image resolution\quad 1080 $\times$ 1920 px\quad 3008 $\times$ 3970 px

Aspect ratio\quad 0.562\quad 0.760

Class distribution\quad Uniform (246–251 images per class, $\pm 3$)\quad Diverse (min 299, max 581 per class)

Shooting conditions\quad Inside a park, 3–5 m distance; various seasons and times; backgrounds include playgrounds and signboards\quad Urban streets, 0.1–3 m distance; various seasons and cities; backgrounds include people, vehicles, and buildings

Capture devices\quad iPhone 15 Pro, Pixel 7 Pro\quad Multiple mobile devices

